\section{מבוא ושאלות המחקר}

שוק ההון אינו מתנהג תמיד באותו אופן. יש תקופות שבהן התשואות קטנות והתנודתיות נמוכה, תקופות שבהן השוק נע במהירות ובחוסר יציבות, ולעיתים גם תקופות של ירידות חדות ולחץ. לכן סביר לחשוב על סדרת המחירים לא כתהליך סטטיסטי יחיד וקבוע, אלא כתהליך שעובר בין כמה ``מצבים'' שונים לאורך הזמן.

מודלים מסוג \tech{regime switching} מתארים את הרעיון הזה באמצעות משתנה חבוי, שאינו נצפה ישירות אך משפיע על ההתפלגות של הנתונים שנמדדים בפועל \cite{hamilton1989new}. ב־\tech{Hidden Markov Model (HMM)}, המצב החבוי משתנה לאורך הזמן לפי הסתברויות מעבר, וכל מצב מייצר תצפיות בעלות מאפיינים שונים \cite{rabiner1989tutorial,murphy2012machine}. בהקשר פיננסי, המשמעות היא שהמודל יכול לזהות למשל תקופה שקטה לעומת תקופה תנודתית, גם אם לא סיפקנו לו מראש תוויות בשם ``שקט'' או ``לחוץ''.

השאלה המקורית של הפרויקט הייתה חיזויית: האם \tech{Gaussian HMM} יכול לזהות מצבי שוק סמויים, ולנצל אותם כדי לשפר את חיזוי כיוון התשואה ביום הבא. קיימות עבודות שמיישמות HMM לנתוני מניות ולבעיות חיזוי דומות \cite{catello2023hmm}. עם זאת, לאחר הריצות הראשונות התברר שהסתכלות רק על דיוק החיזוי מפספסת חלק מהמידע שהמודל מספק. גם כאשר ה־DPA אינו גבוה, ייתכן שה־HMM עדיין מזהה מבנה משמעותי של תנודתיות, סיכון והתמדה לאורך זמן.

לכן העבודה הסופית בוחנת שלוש שאלות מחקר נפרדות:

\begin{enumerate}
    \item \textbf{יכולת חיזוי.} האם HMM משפר את חיזוי הכיוון ביום הבא ביחס לכללי ייחוס פשוטים ולמודלי \tech{Supervised Learning} שמקבלים בדיוק את אותם מאפיינים?
    \item \textbf{איכות הייצוג של מצבי השוק.} האם המצבים החבויים נבדלים זה מזה באופן ברור בתשואה, תנודתיות, טווח יומי, התנהגות נפח המסחר, \tech{drawdown}, שכיחות ומשך שהייה?
    \item \textbf{מידע נוסף על סביבת השוק.} האם זהות המצב ורמת אי־הוודאות של המודל קשורות למעברי משטר, למדד \asset{VIX} ולשינויים במבנה הקורלציה בין נכסים?
\end{enumerate}

ההבחנה בין חיזוי לבין ייצוג היא מרכזית בדוח. חיזוי שואל מה יקרה ביום הבא; ייצוג שואל באיזו סביבה סטטיסטית השוק נמצא כעת. אלה שתי יכולות שונות. מטרת העבודה אינה למצוא בכל מחיר ``מודל מנצח'', אלא להבין איזה מידע ה־HMM מצליח לחלץ מהנתונים, היכן הוא אינו מספק יתרון, ומה ניתן להסיק מהתוצאות באופן זהיר.
